\section{认知重塑与行动指南}

全书从一个落差开始：方案以推理速度出现，现实通过多只时钟改变。把这条转换链看清，才能知道下一次提速究竟改善了什么。

\subsection{走出图灵陷阱：能力、价值与资格分别判断}

行为测试可以检验指定接口上的能力。它的结果属于测试条件，不能替代对持续运行、权利与责任的全部判断。一个系统能写出答案，能稳定提供服务，能代表他人作决定，是三种不同要求。

人的价值也不等于某项产出的稀缺程度。能力扩展人的选择，报酬受需求、产权与议价影响，尊严则不能随一种交付物的市场价格一起折算。生成变便宜以后，社会仍然可以把新增能力用于健康、教育、照护、关系、探索与公共参与。

这里的“负熵体验者”强调两件事：人依赖实际条件维持生命，也在有限生命中体验、选择和建立关系。热力学说明维持的物理条件，伦理与制度说明怎样的生命值得被保护，二者各自提供理由。

价值重新落在哪里，需要观察新的接口。候选生成可能让位于审查；自动检查成熟以后，规范与任务定义又可能成为重点；部署扩大以后，维护、救济与分配可能更加紧迫。接口会移动，职业清单也会过时。

每个位置都要重新问：它建立了什么可依赖结果，凭据在哪里，换一种安排能否更好地完成同一功能，谁能进入这套安排。社会必要性不会自动兑现为工作机会和高报酬，培训、支付与权力配置必须同时处理。

\subsection{价值流向：跟踪稀缺的形成方式}

稀缺有不同来源。验证需要与任务相称的证据；照护需要时间、技能、可达性，有时也需要特定的人；现场判断需要信息、能力与实际处置权；长期教育需要让学生形成能力的路径。工具可以改善每一项，承担这些功能的结构也要随之更新。

评价不能只看名称。审计可能变成盖章，经验要求可能变成排斥，所谓陪伴也可能只剩服务话术。需要检查的是结果与形成路径：出了错能否发现，依据能否复核，承诺是否兑现，受影响者能否提出异议。

高价同样不是必要性的证明。市场稀缺可能来自真实供给约束，也可能来自垄断和人为门槛；重要工作则可能因为预算不足而长期低薪。价值流向的分析，要与分配分析并排进行。

这条规则可以反复重算：每次成本下降以后，辨认旧凭据失去了什么区分力，新凭据需要通过什么过程形成，以及谁掌握形成过程。技术可以创造实际使用价值，也可以重排价值证明的位置；具体收益由完整结果判断。

\subsection{降阻之前的三问}

\textbf{被降掉的是什么？}先固定任务、质量、基准、时间与系统边界。起草变快是一项收益，获得可靠决定变快是另一项，需要分别取证。

\textbf{其他环节怎样变化？}检查复核、返工、能源、设备、支持、风险与恢复。成本可能转移，也可能净下降，还可能出现新的支出。不能预先假定每次降阻都必须在别处原样偿还。

\textbf{谁获得收益，谁承担余账？}把用户、劳动者、运营者、受影响第三方与未来义务列入比较。总成本下降，也可能伴随不均等的收益和负担。

对软件团队，三问把生成连接到审查、回归风险与维护。对临床服务，它把模型表现连接到具体用途、医师义务、患者权利与随访。对运营车队，它把车辆效率连接到空驶、覆盖、设施与事故处置。答案应当具体到能够改变方案。

审查过载，就比较候选筛选和处理能力；故障无法收束，就缩小作用范围或重做恢复；责任没有落点，就先建立可追问与可补偿的安排。表格只记录这些决定，实质检查仍须有人有能力、有时间并有权执行。

\subsection{保留能生产能力的接口}

个体可以保留一部分直接练习，使自己仍能理解、检查与修订所依赖的结果。学习语言时先写再改，维护代码时实际阅读和测试，学习判断时接触结果并解释失败。好的工具可以让练习更有效。

保留多少由任务决定。某些能力需要持续暴露，某些偶尔练习即可维持；有些旧步骤已经不再产生学习，应当删除。目的在能力，而不在展示手工、吃力或不用工具的身份。

个人条件也不同。时间、残障、照护、收入与疲劳改变可选安排，辅助工具往往正是在高障碍处扩大自主。组织不能一面只考核即时产出，一面要求劳动者在私人时间无偿保存所有长期能力。

组织需要维护的是复核与重建过程的能力：相称的记录、胜任的审查、独立质疑，以及真正能执行的中断和纠正权。这些功能需要角色、预算与权限，不能用一句“谨慎使用”替代。

\subsection{让承担者保有有效位置}

受影响者需要知情、参与、代表与救济。有些决定可由本人同意，有些需要专业职责、公共程序或对脆弱群体的保护。决策权要与能力和控制相称，承担损失的人则需要实际可用的质疑和补偿路径。

下一代工具会改变许多具体答案。这套方法仍可继续使用：辨认介质，分开约束，追踪行动，检查反馈，核对整本账。能力越强，越需要知道它通过什么路径成为现实。
